\section{Sensitivity bands under political bias}\label{sec:bias}

The bounds in Sections~\ref{sec:bounds}--\ref{sec:contradiction} treat each measured input as honest. We now allow the analyst to dial in distrust of specific sources. Political bias enters the confidence statement through Huber contamination: each source $i$ writes its likelihood as $(1-\varepsilon_i)L_i+\varepsilon_i Q_i$, with $Q_i$ adversarial in some class $\mathcal Q$ \citep{huber1964,huber1981,berger1986,gagnon2023}. The contamination parameter $\varepsilon_i\in[0,1]$ is the analyst's prior probability that source $i$ is producing a politically motivated reading rather than a measurement of the true distribution.

\begin{proposition}[Bias-robust sensitivity band]\label{prop:robust}
Let $\theta$ be a scalar estimand with $1-\alpha$ posterior credible interval $I_0=[L_0,U_0]$ under face-value sources, with the posterior supported on the identified set implied by the sources used. Suppose each source $i$ is independently contaminated with probability $\varepsilon_i$, and removing the sources in a set $T$ leaves an identified set of diameter $R_T$ (write $R_i:=R_{\{i\}}$).
\begin{enumerate}[label=(\roman*),leftmargin=*]
\item \emph{(Patternwise envelope.)} If the realised contamination pattern is $T$, every \emph{support-respecting} $1-\alpha$ credible interval---one whose endpoints lie in the posterior support, e.g.\ the equal-tailed or HPD interval---obtainable by adversarial replacement of the sources in $T$ lies within $[L_0-R_T,\ U_0+R_T]$. (An arbitrary credible interval can be widened without limit while retaining $1-\alpha$ coverage, so no finite envelope constrains that class.)
\item \emph{(Expected displacement.)} Averaging over the contamination pattern, the expected endpoint displacement is at most $\sum_i\varepsilon_i R_i$ plus higher-order terms in $\varepsilon$ involving the multi-removal diameters $R_T$, $|T|\ge 2$; when at most one source can be contaminated, the bound $\sum_i\varepsilon_i R_i$ is exact. The band
\begin{equation}\label{eq:Iepsilon}
I_\varepsilon \;=\; \Bigl[L_0-\textstyle\sum_i\varepsilon_i R_i,\; U_0+\textstyle\sum_i\varepsilon_i R_i\Bigr]
\end{equation}
is therefore a first-order sensitivity band for the interval endpoints---not a uniform envelope over realised contaminated intervals, whose endpoints can move by up to $R_T$ on the (probability $\le\varepsilon_i$) contamination events.
\item \emph{(Quantile displacement.)} If additionally the posterior mixture weight on contaminated components is at most $\tau:=1-\prod_i(1-\varepsilon_i)$---which holds when the component marginal likelihoods are comparable, and must otherwise be checked, since posterior weights are proportional to prior pattern probabilities \emph{times} component marginal likelihoods---and the face-value posterior density exceeds $m>0$ on the segment between the clean and contaminated $u$-quantiles, then
\[
\bigl|\hat Q_\varepsilon(u)-\hat Q_0(u)\bigr|
   \;\le\; m^{-1}\tau,
\qquad u\in\{\alpha/2,\,1-\alpha/2\}.
\]
\end{enumerate}
\end{proposition}
Proof in Appendix~\ref{app:proofs}.

\paragraph{Practice.} Set $\varepsilon_i\approx 0.05$ for widely-trusted demographic sources (Palestinian Central Bureau of Statistics, health-ministry demographic records consistent with independent surveys) and $\varepsilon_i\approx 1$ (i.e., adversarial) for belligerent claims; \eqref{eq:Iepsilon} then absorbs the assumed political bias automatically. We apply the band to the identified-set endpoints---the degenerate case in which the posterior is supported on the identified set and the credible interval is the set itself. In the Gaza application the removal diameters, each computed as the diameter of the exposure-agnostic identified set with that source removed and the remaining inputs held at their envelopes, are: $R_{\mathrm{demo}}=\nRDemo{}$ (dropping the demographic record leaves only the manpower bound $q\le M/D$); $R_{\mathrm{census}}=\nRCensus{}$ (letting $w$ range over a generous global envelope $[0.70,0.76]$, with $a=1-w$, gives the set $[0,\,q(1;w{=}0.76)]=[0,\nRCensusSet{}]$); and $R_M=\nRM{}$ (dropping the manpower source frees $f$ up to its logical cap $f\le a$, at which the set reaches the all-male ceiling $1-\omega$). The global envelope $[0.70,0.76]$ is a background constraint---a generous envelope for the demographic structure of young, high-fertility populations of the kind at issue here, justified independently of the census being removed---so it survives removal of the census source, and the removal diameter is computed under it. Then $\varepsilon=0.05$ inflates the exposure-agnostic upper bound by $0.05\,(\nRDemo{}+\nRCensus{}+\nRM{})\approx \nHuberPP{}$ percentage points---from $\nQOne{}\%$ to $\nHuberInflated{}\%$---which still rejects the upper end of the disputed claim.
